\section{Boundaries of the Observation}
\label{sec:analysis}

\subsection{Descriptive variation across MMLU groups}
Across the standard broad MMLU groups, keep14 remains below the intact base and
ShortGPT remains above keep14. The realized keep14 scores range from .294 on
STEM to .377 on Other; ShortGPT improves each group but does not close the
base gap. These are descriptive fixed-checkpoint summaries. We did not
predeclare a family of domain tests, adjust for the 57 subjects examined, or
estimate variation across training runs. We therefore do not infer
domain-specific recovery effects, knowledge types, or anatomical boundaries
from these differences. Complete group and subject tables are in
Appendix~\ref{app:mmlu-subjects}.

\subsection{The shallow observations are not a depth law}
The retained shallow rows stop at 121k (keep8), 83.5k (keep10), and 124k
(keep12), whereas keep14 reaches 200k. The checkpoints were retained after
knowledge-sensitive metrics appeared stable, without a registered stopping
rule, and PPL was still decreasing. Equal steps would also imply unequal FLOPs
across depths. Consequently, the ladder is useful only as an inventory of
available operating points. It cannot support a converged architectural
threshold, a depth-dependent recovery rate, or a compute-normalized ordering.

Within keep8, several likelihood-style tasks and PPL improve over the observed
interval while aggregate MMLU remains near chance. Within keep14, MMLU rises
modestly late in training. These two literal traces differ, but their starting
quality, budget, and stochastic realization are unmatched. No claim about a
general dynamic follows.

\subsection{Scope checks do not constitute replications}
A more compressed OLMo-2-1B prefix+fresh-tail run shows falling PPL with
near-chance MMLU, and an existing Qwen3-8B endpoint is also near chance after
200k. The 1B arm changes scale and retained fraction; the Qwen arm additionally
changes architecture, corpus, and available evaluations. We include them as
directional context only, not as matched replications or evidence of
cross-family universality (Appendix Table~\ref{tab:app-crossfamily}).

\subsection{What remains identified}
The evidence identifies a measurement limitation: for the reported OLMo
constructions, budgets, and interfaces, in-domain perplexity alone cannot
certify the measured knowledge-sensitive evaluations. The operating points
also bound three complete explanations: short-horizon corpus shift, an
answer-letter-only artifact, and nominal depth alone. They do not identify the
responsible structural factor, estimate a causal treatment effect, establish
what happens after the observed horizon, or locate any capability in the
network.
